\documentclass[10pt,a4paper]{amsart}
\usepackage[margin=19mm]{geometry}
\usepackage{amsmath,amssymb,mathtools}
\usepackage[scr=rsfs,cal=euler]{mathalfa}
\usepackage{xfrac}
\usepackage[unicode,hidelinks,bookmarks=false]{hyperref}
\newcommand{\ie}{i.e.\ }
\newcommand{\cf}{cf.\ }
\newcommand{\resp}{respectively\ }
\newcommand{\Subsection}{\subsection}
\providecommand{\nobreakdash}{\nobreak\hbox{-}}
\setlength{\emergencystretch}{2em}
\multlinegap0pt
\usepackage{mathtools}
\usepackage{amssymb}
\usepackage[shortlabels]{enumitem}
\usepackage{xcolor}
\newtheorem{theorem}{Theorem}
\newtheorem{proposition}{Proposition}
\newtheorem{lemma}{Lemma}
\newcommand{\RR}{\mathbb R}
\newcommand{\epsov}{\varepsilon}
\title{The Ferrers bound for spanning trees in bipartite graphs}
\author{Boon Suan Ho}
\date{Source: arXiv:2603.17997v1 (CC BY 4.0)}
\begin{document}
\maketitle

\noindent\emph{Source and licence.} The Ferrers bound for spanning trees in bipartite graphs. Version arXiv:2603.17997v1 (CC BY 4.0), distributed by arXiv under the Creative Commons Attribution 4.0 International licence, http://creativecommons.org/licenses/by/4.0/, as recorded in the arXiv licence metadata for this exact version and shown on the abstract page https://arxiv.org/abs/2603.17997v1. Full licence notice: \texttt{licenses/arxiv-2603.17997-CC-BY-4.0.txt}.\par\smallskip

\noindent\textit{Editorial scope.} Counted unit: the article's theorem thm:main (source0.tex lines 203--210). The article's complete original proof of it is delivered with it (source0.tex lines 719--831, one \texttt{proof} environment of 113 lines, which the article places away from the statement). No mathematics is written, repaired or re-derived: the statement and the proof are the article's own lines, in the article's own notation, and no retained line is substituted or re-flowed.  Retained uncounted as the source-local context the proof needs: lines 274--282; lines 397--404; lines 431--440; lines 437--439; lines 535--542; lines 591--600; lines 657--684; lines 671--678.  Nothing was omitted from any retained span, so the package carries no omission record.  No retained line contains a citation, so the wrapper carries no bibliography and no bibliography entry.  No retained line contains a figure environment, an image-inclusion command or drawing code, so no image is delivered and the package declares no asset dependency.  Editorial formatting only, in addition: an AMS \texttt{amsart} preamble that restates the article's own declarations and macros used by the retained lines, namely the environments \texttt{theorem}, \texttt{proposition}, \texttt{lemma} on separate counters, together with the standard packages those lines need.  The retained environments print in this document as Theorem 1, Proposition 1, Lemma 1 in order of appearance, whereas the article prints the same items with the numbering of its own full document.  The complete native source of the article as distributed in the arXiv e-print for this exact version is bundled unchanged as \texttt{sources/196665/source0.tex}.
\begin{proposition}\label{prop:strict-schur}
If $x,y\in(0,\infty)^m$ and $x$ majorizes $y$, then for every strictly 
concave function $\phi\colon(0,\infty)\to\RR$,
\[
\sum_{i=1}^m \phi(x_i)\le \sum_{i=1}^m \phi(y_i).
\]
If, in addition, $x$ and $y$ are both written in weakly decreasing order, 
then equality holds if and only if $x_i=y_i$ for all $i$.
\end{proposition}

\begin{lemma}\label{lem:kernel}
The matrix $L_X$ is symmetric positive semidefinite, satisfies 
$L_X\mathbf1=0$, and has kernel exactly $\langle \mathbf1\rangle$. 
In particular, its eigenvalues are
\[
0=\mu_1<\mu_2\le \cdots \le \mu_m.
\]
\end{lemma}

\begin{proposition}\label{prop:reduction}
Let $M\coloneqq L_X+\frac{n}{m}J$, where $J$ is the $m\times m$ all-ones matrix. Then
\begin{equation}\label{eq:tau-detM}
\tau(G)=\frac{\prod_{j=1}^n b_j}{mn}\det M.
\end{equation}
Consequently, the inequality in Theorem~\ref{thm:main} is equivalent to
\begin{equation}\label{eq:det-bound}
\det M\le \prod_{i=1}^m a_i.
\end{equation}
\end{proposition}

\begin{equation}
\det M\le \prod_{i=1}^m a_i.
\end{equation}

\begin{lemma}\label{lem:Q-proj}
For every nonempty $T\subseteq X$, the matrix $Q_T$ is the 
orthogonal projection onto
\[
H_T\oplus \langle \mathbf1\rangle.
\]
In particular, $Q_T$ has rank $|T|$.
\end{lemma}

\begin{lemma}\label{lem:overlap}
For nonempty subsets $I,T\subseteq X$,
\begin{equation}\label{eq:overlap}
\operatorname{tr}(Q_IQ_T)
=
|I\cap T|+\frac{|I\setminus T||T\setminus I|}{|I||T|}
\ge |I\cap T|.
\end{equation}
Equality holds if and only if $I\subseteq T$ or $T\subseteq I$.
\end{lemma}

\begin{proposition}\label{prop:majorization}
Let
\[
\lambda_1\ge \lambda_2\ge \cdots \ge \lambda_m
\]
be the eigenvalues of $M$, and relabel the vertices of $X$ so that
\[
a_1\ge a_2\ge \cdots \ge a_m.
\]
For each $k\in\{1,\dots,m\}$, write
\[
[k]\coloneqq \{x_1,\dots,x_k\}.
\]
Then
\begin{equation}\label{eq:partial-sum-ineq}
\sum_{i=1}^k \lambda_i
\ge
\sum_{i=1}^k a_i+\sum_{j=1}^n \epsov([k],T_j)
\ge
\sum_{i=1}^k a_i
\qquad (1\le k<m),
\end{equation}
and
\[
\sum_{i=1}^m \lambda_i=\sum_{i=1}^m a_i.
\]
Consequently, $(\lambda_1,\dots,\lambda_m)$ majorizes $(a_1,\dots,a_m)$.
\end{proposition}

\begin{equation}
\sum_{i=1}^k \lambda_i
\ge
\sum_{i=1}^k a_i+\sum_{j=1}^n \epsov([k],T_j)
\ge
\sum_{i=1}^k a_i
\qquad (1\le k<m),
\end{equation}

\begin{theorem}\label{thm:main}
Let $G=(X\sqcup Y,E)$ be a connected bipartite graph with $|X|=m$ and $|Y|=n$. 
Then
\[
\tau(G)\le \frac{1}{mn}\prod_{v\in V(G)}\deg(v).
\]
Moreover, equality holds if and only if $G$ is a Ferrers graph.
\end{theorem}

\begin{proof}[Proof of Theorem~\ref{thm:main}]
Let
\[
\lambda_1\ge \lambda_2\ge \cdots \ge \lambda_m
\]
be the eigenvalues of $M$, and relabel the vertices of $X$ so that
\[
a_1\ge a_2\ge \cdots \ge a_m.
\]
For each $k\in\{1,\dots,m\}$, write
\[
[k]\coloneqq \{x_1,\dots,x_k\}.
\]
By Proposition~\ref{prop:majorization}, 
the vector $(\lambda_1,\dots,\lambda_m)$ majorizes 
the degree vector $(a_1,\dots,a_m)$.

By Lemma~\ref{lem:kernel}, $L_X$ has eigenvalues $0,\mu_2,\dots,\mu_m$, 
and since $J$ acts by $m$ on $\langle\mathbf1\rangle$ and by $0$ on 
$\mathbf1^\perp$, $M$ has eigenvalues $n,\mu_2,\dots,\mu_m$, and is 
thus positive definite, so we have $\lambda_i>0$ for every $i$. 
Because $G$ is connected, every vertex degree is positive, so $a_i>0$ 
for every $i$ as well. Thus both vectors lie in $(0,\infty)^m$, and 
Proposition~\ref{prop:strict-schur} applied with the strictly concave 
function $\phi(t)=\log t$ gives
\[
\sum_{i=1}^m \log \lambda_i\le \sum_{i=1}^m \log a_i.
\]
Equivalently,
\[
\det M=\prod_{i=1}^m \lambda_i\le \prod_{i=1}^m a_i.
\]
By Proposition~\ref{prop:reduction},
\[
\tau(G)=\frac{\prod_{j=1}^n b_j}{mn}\det M
\le
\frac{1}{mn}\Bigl(\prod_{i=1}^m a_i\Bigr)\Bigl(\prod_{j=1}^n b_j\Bigr)
=
\frac{D(G)}{mn},
\]
which proves the inequality.

Now suppose equality holds in Theorem~\ref{thm:main}. 
Then equality holds in \eqref{eq:det-bound}, so
\[
\sum_{i=1}^m \log \lambda_i=\sum_{i=1}^m \log a_i.
\]
We already know from Proposition~\ref{prop:majorization} that 
$(\lambda_1,\dots,\lambda_m)$ majorizes $(a_1,\dots,a_m)$, and both 
vectors were arranged in weakly decreasing order at the start of the proof. 
Therefore the equality statement in Proposition~\ref{prop:strict-schur} 
yields coordinatewise equality,
\[
\lambda_i=a_i\qquad (1\le i\le m).
\]
Returning to \eqref{eq:partial-sum-ineq}, we get
\[
\sum_{j=1}^n \epsov([k],T_j)=0
\qquad (1\le k<m).
\]
Each summand is nonnegative, so every summand vanishes:
\[
\epsov([k],T_j)=0
\qquad\text{for all }j\text{ and }1\le k<m.
\]
By Lemma~\ref{lem:overlap}, each $T_j$ is comparable by inclusion with every 
initial segment $[k]$.

We now show that this forces $T_j$ itself to be an initial segment. Fix $j$ and let
\[
t\coloneqq \max\{i:\ x_i\in T_j\}.
\]
Then $T_j\subseteq [t]$. If $t=1$, then $T_j\neq\varnothing$ implies $T_j=[1]$. 
Assume from now on that $t\ge 2$. Since $x_t\in T_j$, the inclusion 
$T_j\subseteq [t-1]$ is impossible. Comparability of $T_j$ with $[t-1]$ 
therefore gives
\[
[t-1]\subseteq T_j.
\]
Together with $x_t\in T_j$, this yields $T_j=[t]$. Thus every neighborhood of 
a $Y$-vertex is an initial segment of $X$. Reordering the vertices of $Y$ so 
that these initial segments appear in decreasing size, 
we conclude that $G$ is Ferrers.

Conversely, suppose $G$ is Ferrers. Then after ordering $X$ and $Y$ we have
\[
T_j=[t_j]\qquad\text{with}\qquad t_1\ge\cdots\ge t_n\ge 1.
\]
The subspaces
\[
\operatorname{im}(Q_{[1]})\subset \operatorname{im}(Q_{[2]})
\subset \cdots \subset \operatorname{im}(Q_{[m]})=\RR^m
\]
form a complete flag, and $\dim \operatorname{im}(Q_{[t]})=t$ by 
Lemma~\ref{lem:Q-proj}. Choose an orthonormal basis $u_1,\dots,u_m$ adapted 
to this flag, meaning that $u_1,\dots,u_t$ span $\operatorname{im}(Q_{[t]})$ 
for every~$t$. In this basis each $Q_{[t]}$ is diagonal with exactly the 
first~$t$ diagonal entries equal to~$1$. Hence
\[
M=\sum_{j=1}^n Q_{T_j}=\sum_{j=1}^n Q_{[t_j]}
\]
is diagonal, with diagonal entries
\[
\#\{j:\ t_j\ge 1\},\ \#\{j:\ t_j\ge 2\},\ \dots,\ \#\{j:\ t_j\ge m\}.
\]
But $\#\{j: t_j\ge r\}$ is exactly the degree of $x_r$, since $x_r$ is 
adjacent precisely to those~$y_j$ with $t_j\ge r$. Therefore the diagonal 
entries are $a_1,\dots,a_m$, so
\[
\det M=\prod_{i=1}^m a_i.
\]
Proposition~\ref{prop:reduction} now gives equality in Theorem~\ref{thm:main}.
\end{proof}

\end{document}
